\documentclass[11pt]{article}
\usepackage[margin=1in]{geometry}
\usepackage{amsmath,amssymb}
\usepackage{booktabs}
\usepackage{multirow}
\usepackage{graphicx}
\usepackage{xcolor}
\usepackage{hyperref}
\usepackage[numbers]{natbib}

\title{Credit Assignment Granularity in Multi-Turn Agent RL:\\An Empirical Taxonomy of Failure Modes}

\author{Anonymous}
\date{}

\begin{document}
\maketitle

%% ============================================================
%% ABSTRACT
%% ============================================================
\begin{abstract}
Reinforcement learning (RL) for large language model (LLM) agents in multi-turn interactive environments requires assigning credit across long action sequences.
Community frameworks default to \emph{step-independent} training, which splits trajectories into fixed-length windows optimized separately.
We conduct a systematic empirical study of how credit assignment granularity affects training outcomes in ALFWorld, a household task benchmark with 15--45 step episodes.
Our experiments reveal three findings.
First, step-independent training exhibits a stochastic performance cliff: at small context windows it frequently collapses to near-zero success, but the cliff location varies across random seeds (mean 1.0\% at $h{=}2$, std 1.8pp; mean 11.5\% at $h{=}5$, std 10.4pp), precluding a deterministic threshold claim.
Second, episode-level Group Relative Policy Optimization (GRPO), which treats entire trajectories as single optimization units, reaches 23.8\% mean success but plateaus by step 10--15 and is outperformed by Generalized Advantage Estimation with $\lambda{=}0.95$ (43.8\%, $+$20.0pp), demonstrating that finer per-token credit is necessary for strong performance.
Third, selective self-verification trained on top of the episode-level agent fails: the learned policy degenerates to near-always-verify (88\% verification rate, 12\% cost reduction) and is Pareto-dominated by naive baselines.
These results map the failure boundary of coarse credit assignment and identify base agent competence ($<$30\% success) as a prerequisite bottleneck for selective verification.
\end{abstract}

%% ============================================================
%% 1  INTRODUCTION
%% ============================================================
\section{Introduction}

Training LLM-based agents with RL to solve multi-turn interactive tasks has attracted growing interest.
A central design choice in this setting is \emph{credit assignment granularity}: how to distribute a sparse episode-level reward signal across the many tokens and actions that compose a trajectory.
Recent work pursues increasingly fine-grained credit, including turn-level value functions, step-level process reward models, and hierarchical planners.
The implicit consensus is that finer credit assignment is uniformly better.

This paper tests that assumption empirically.
We study three levels of granularity in ALFWorld, a benchmark where agents execute 15--45 step household tasks: (1) \textbf{step-independent} training, which partitions trajectories into fixed-length windows of $h$ steps and optimizes each window separately using discounted returns; (2) \textbf{episode-level} GRPO, which treats the full trajectory as a single sample with a shared advantage; and (3) \textbf{GAE-based} training, which uses a learned value function to compute per-token temporal-difference advantages.

Step-independent training is the default in community frameworks such as RAGEN, yet its failure modes in long-horizon agent tasks are poorly characterized.
We hypothesize two mechanisms behind its failures: \emph{policy information loss} (short windows prevent learning cross-step strategies) and \emph{advantage sign misattribution} (per-window advantage signs are dominated by episode outcome rather than local action quality).

Our experiments yield a sobering picture.
Step-independent training does fail at small $h$, but the failure is stochastic rather than catastrophic: at $h{=}5$, one of three seeds collapses while two reach 15--19\% success.
Episode-level GRPO avoids this collapse but plateaus at 23.8\%, far below GAE's 43.8\%, confirming that per-token credit provides a substantial advantage.
Selective verification, our originally proposed application, fails entirely because the base agent is too weak ($<$30\% success) for verification to carry meaningful signal.

We structure these findings as an empirical taxonomy of credit assignment failure modes, organized by three axes: granularity level, agent competence, and task horizon.
Our contributions are:
\begin{itemize}
\item A multi-seed characterization of the step-independent performance cliff, showing it is seed-dependent (collapse rate 67\% at $h{=}2$, 33\% at $h{=}5$, 0\% at $h{\geq}10$) rather than a deterministic threshold.
\item Mechanistic evidence that advantage sign misattribution exists but is weaker than theoretically predicted (sign-outcome correlation 0.72 vs. predicted $>$0.95; 18\% mixed-sign episodes vs. predicted 0\%).
\item A negative result demonstrating that selective verification requires a base agent competence threshold: at $<$30\% success, learned verification policies degenerate and are dominated by random baselines.
\item Cross-environment (ScienceWorld) and cross-scale (7B) experiments showing the step-independent disadvantage is environment-specific and diminishes with model scale.
\end{itemize}


%% ============================================================
%% 2  BACKGROUND AND RELATED WORK
%% ============================================================
\section{Background and Related Work}

\paragraph{Multi-turn agent RL.}
Training LLM agents for sequential decision-making environments has been explored through several frameworks.
RAGEN implements GRPO for multi-turn agents with a step-independent training mode as its default.
ArCher establishes hierarchical RL for multi-turn agents.
HiPER achieves 97.4\% on ALFWorld using hierarchical plan-execute RL at 7B scale.
GAGPO proposes grouped value proxies aligned to action steps, reaching 93.5\% on ALFWorld at 1.5B.
These systems uniformly pursue finer credit assignment for higher performance.

\paragraph{Credit assignment in RL for LLMs.}
Turn-PPO identifies the ``uniform advantage'' problem where all turns in a trajectory receive identical advantage estimates, and proposes turn-level value functions.
MT-GRPO designs turn-level rewards for search-based QA tasks.
iStar uses implicit process reward models with dual advantages.
A recent survey by Zhang et al.\ catalogues 47 credit assignment methods but does not analyze when coarser methods suffice or fail.
Our work fills this gap by empirically characterizing failure boundaries.

\paragraph{Self-verification for agents.}
Self-verification, where an agent re-evaluates its action before committing, improves performance in interactive tasks.
Self-Guide uses always-on verification with expert-guided rewards, reaching 72.7\% on ALFWorld at 1.7B scale.
ARES trains a supervised router to selectively trigger verification, achieving 52.7\% token reduction.
ReflAct uses prompting-based reflection with GPT-4.
Our work attempted to train a selective verification policy via RL, but the negative result below reveals a prerequisite that prior work does not discuss: base agent competence must exceed a minimum threshold for verification to be learnable.


%% ============================================================
%% 3  METHOD
%% ============================================================
\section{Method}

\subsection{Problem Setting}

We consider an LLM agent interacting with a text-based environment over $T$ steps.
At each step $t$, the agent observes state $s_t$ and generates action $a_t$.
A binary reward $R \in \{0, 1\}$ is given at episode end based on task completion.
The agent optionally performs self-verification: after generating $a_t$, it re-reads the context and may revise its action to $a_t'$ before execution.

\subsection{Credit Assignment Variants}

\paragraph{Step-independent GRPO.}
The trajectory is partitioned into non-overlapping windows of $h$ steps.
Each window $w_i = (s_{ih}, a_{ih}, \ldots, s_{(i+1)h-1}, a_{(i+1)h-1})$ receives a $\gamma$-discounted return computed from the episode reward.
GRPO computes group-relative advantages within a batch: $\hat{A}_i = (R_i - \mu_{\text{group}}) / \sigma_{\text{group}}$.
In step-independent mode, each window is optimized as an independent sample.

\paragraph{Episode-level GRPO.}
The entire trajectory of $T$ steps is treated as a single sample.
All tokens share one advantage value $\hat{A} = (R - \mu_{\text{group}}) / \sigma_{\text{group}}$.
This eliminates per-window return computation and the associated $\gamma$-discounting, but provides no per-token credit differentiation.

\paragraph{GAE baseline.}
Generalized Advantage Estimation uses a learned value function $V_\phi(s_t)$ to compute per-token advantages:
$\hat{A}_t^{\text{GAE}} = \sum_{l=0}^{T-t} (\gamma \lambda)^l \delta_{t+l}$,
where $\delta_t = r_t + \gamma V_\phi(s_{t+1}) - V_\phi(s_t)$.
This provides the finest credit granularity among our methods.

\subsection{Selective Verification (B-IAVR)}

We train a binary verification policy $\pi_\theta(v_t | s_t, a_t)$, $v_t \in \{\text{verify}, \text{skip}\}$, using GRPO on top of a frozen base agent.
The reward signal for the verification policy is based on \emph{action change}: if verification changes the agent's action ($a_t' \neq a_t$), the verification is deemed informative.
The composite reward combines task success and action-change informativeness.

\subsection{Hypothesized Failure Mechanisms}

We hypothesize two mechanisms for step-independent failure in long-horizon tasks:

\paragraph{Policy information loss.}
A window of $h$ steps provides insufficient context for cross-step planning.
For tasks requiring 30--45 steps, a window of $h{=}5$ captures only 11--17\% of the trajectory.

\paragraph{Advantage sign misattribution.}
In step-independent training, the advantage sign for window $w_i$ is determined primarily by the episode outcome rather than the quality of actions within $w_i$.
In a failed episode, a locally good action receives negative advantage; in a successful episode, a locally poor action receives positive advantage.
The GRPO group normalization may introduce partial local ranking signal, weakening this effect.


%% ============================================================
%% 4  EXPERIMENTS
%% ============================================================
\section{Experiments}

All experiments use Qwen3-1.7B as the base model trained with GRPO on ALFWorld, unless otherwise noted.
Training runs use 15 gradient steps with batch size 8 and 4 rollouts per prompt.
We report validation success rate on a held-out set of 32 tasks.

\subsection{Context Window Ablation (h-Sweep)}

We sweep $h \in \{2, 5, 10, 20, \text{full}\}$ with 3 random seeds each to characterize how the step-independent performance cliff depends on window size.

\begin{table}[h]
\centering
\caption{Step-independent GRPO performance by context window size $h$ (3 seeds, 15 steps).}
\label{tab:hsweep}
\begin{tabular}{lcccc}
\toprule
$h$ & Success (mean) & Std & Collapse Rate & Convergence Step \\
\midrule
2    & 1.0\%  & 1.8pp & 67\% & 3--4 \\
5    & 11.5\% & 10.4pp & 33\% & 4--6 \\
10   & 14.6\% & 6.3pp  & 0\%  & N/A \\
20   & 19.8\% & 4.8pp  & 0\%  & N/A \\
full & 25.0\% & 3.1pp  & 0\%  & N/A \\
\bottomrule
\end{tabular}
\end{table}

Table~\ref{tab:hsweep} reveals a monotonic improvement with increasing $h$, but the transition from failure to success is not a sharp cliff.
At $h{=}2$, two of three seeds collapse to 0\% while one achieves 3.1\%.
At $h{=}5$, the picture is more complex: seed 2 collapses to 0\%, but seeds 1 and 3 reach 15.6\% and 18.8\% respectively, yielding a standard deviation of 10.4pp.
This high variance means the ``cliff location'' is a random variable, not a structural property of the algorithm.
At $h{\geq}10$, no seeds collapse, but performance remains well below episode-level (25.0\%).
The cliff threshold range is seed-dependent, falling between $h{=}3$ and $h{=}8$.

\subsection{Multi-Seed Episode-Level Replication}

We run 5 seeds of episode-level GRPO to establish its mean and variance.

\begin{table}[h]
\centering
\caption{Episode-level GRPO across 5 seeds. All seeds plateau by step 10--15.}
\label{tab:multiseed}
\begin{tabular}{lccccc}
\toprule
Seed & Step 5 & Step 10 & Step 15 & Step 50 & Plateau Onset \\
\midrule
1 & 25.0\% & 25.0\% & 28.1\% & 28.1\% & step 10 \\
2 & 18.8\% & 18.8\% & 18.8\% & 21.9\% & step 5 \\
3 & 25.0\% & 31.3\% & 31.3\% & 31.3\% & step 10 \\
4 & 12.5\% & 15.6\% & 15.6\% & 15.6\% & step 5 \\
5 & 18.8\% & 25.0\% & 25.0\% & 28.1\% & step 10 \\
\midrule
Mean & 20.0\% & 23.1\% & 23.8\% & 25.0\% & -- \\
Std  & 5.1pp  & 6.1pp  & 6.2pp  & 6.3pp  & -- \\
\bottomrule
\end{tabular}
\end{table}

The mean success rate is 23.8\% $\pm$ 6.2pp at step 15 (Table~\ref{tab:multiseed}).
All seeds plateau between steps 5 and 10, and an additional 35 training steps (to step 50) yields only 1.2pp improvement (25.0\%).
This plateau is consistent with an information bottleneck: episode-level credit provides correct advantage direction but no per-token magnitude differentiation.
The 95\% confidence interval [11.3\%, 36.3\%] overlaps with step-independent $h{=}10$ (14.6\%), weakening the statistical separability of these methods.

\subsection{GAE Baseline Comparison}

We compare episode-level GRPO against GAE with two configurations: $\lambda{=}0.95, \gamma{=}1.0$ (standard) and $\lambda{=}1.0, \gamma{=}1.0$ (Monte Carlo-like with critic baseline).

\begin{table}[h]
\centering
\caption{Credit assignment method comparison (3 seeds, 15 steps).}
\label{tab:gae}
\begin{tabular}{lccccc}
\toprule
Method & Success & Std & $\Delta$ vs. Episode & Valid Actions & Improving? \\
\midrule
Episode GRPO    & 23.8\% & 6.2pp & --      & 55\% & No (plateau) \\
GAE($\lambda{=}1.0$)  & 37.5\% & 5.5pp & +13.7pp & 72\% & Marginal \\
GAE($\lambda{=}0.95$) & 43.8\% & 4.7pp & +20.0pp & 78\% & Yes \\
\bottomrule
\end{tabular}
\end{table}

GAE($\lambda{=}0.95$) outperforms episode-level GRPO by 20.0pp (Table~\ref{tab:gae}), a substantial gap.
This demonstrates that per-token temporal-difference credit is important: knowing the episode outcome is insufficient; the agent needs to know \emph{which tokens contributed} to success or failure.
GAE also achieves higher action validity (78\% vs. 55\%) and continues improving at step 15, unlike episode-level GRPO which has plateaued.
The GAE($\lambda{=}1.0$) variant, equivalent to Monte Carlo returns with a critic baseline, achieves 37.5\%, still 13.7pp above episode-level GRPO.
These results confirm that finer credit assignment is necessary for strong performance in long-horizon agent tasks.

\subsection{Credit Assignment Mechanism Verification}

We empirically test the advantage sign misattribution hypothesis by logging per-window advantages and correlating them with episode outcomes and local action quality.

\begin{table}[h]
\centering
\caption{Advantage sign analysis: theoretical predictions vs. observed values.}
\label{tab:mechanism}
\begin{tabular}{lccc}
\toprule
Metric & Predicted & Observed & Interpretation \\
\midrule
Sign-outcome correlation   & $>$0.95 & 0.72 & Exists but not dominant \\
Sign-local quality corr.   & $<$0.1  & 0.28 & Non-trivial local signal \\
$\gamma^k$ decay $R^2$     & $>$0.8  & 0.45 & Weak fit \\
Mixed-sign episodes         & 0\%     & 18\% & Theory violation \\
Episode-level adv.\ variance & 0       & 0    & Confirmed (trivially) \\
\bottomrule
\end{tabular}
\end{table}

Table~\ref{tab:mechanism} shows that sign misattribution is a real phenomenon but weaker than predicted.
The sign-outcome correlation of 0.72 confirms that advantage signs are biased toward episode outcomes, but 28\% of the signal reflects local action quality.
The 18\% mixed-sign episodes violate the pure misattribution theory, which predicts 0\%: GRPO's group normalization introduces within-group ranking that partially preserves local information.
The $\gamma^k$ magnitude decay shows only moderate fit ($R^2 = 0.45$), meaning advantage magnitudes decay with window position but noisily.
Sign misattribution is \emph{a contributing factor} to step-independent failure, not \emph{the} explanation.

\subsection{Valid Action Bonus Ablation}

We test whether adding a small dense reward for action validity ($+$0.05 per valid action) can break the episode-level plateau.

Episode-level GRPO with the valid action bonus achieves 25.0\% success at step 15 (vs. 23.8\% vanilla), a negligible improvement.
However, action validity improves substantially: from 55\% to 72\% at step 15 and 85\% at step 30.
The agent learns to produce syntactically correct actions but not strategically better ones.
The plateau is not caused by invalid action formatting; it stems from the inability of episode-level credit to differentiate among valid actions.
This confirms that the bottleneck is credit assignment granularity, not reward sparsity.

\subsection{Selective Verification Results}

\subsubsection{Action-Change Statistics}

Before training the selective policy, we analyze whether action change after verification correlates with success in the base agent.
The action change rate is 35\%, with 65\% of verifications being confirmatory.
The correlation between action change and episode success is 0.08, and after controlling for policy entropy, the partial correlation drops to 0.03 ($p = 0.42$, not significant).
Action change in the low-entropy stratum correlates at 0.05, in the high-entropy stratum at 0.11.
The entropy-success correlation is $-$0.15.
These results indicate that action change is a side effect of high policy entropy, not a signal of verification informativeness.

A causal intervention experiment comparing verification-triggered action changes against random changes at the same rate confirms this: verification-triggered changes yield 22\% success vs. random changes at 19\% ($\Delta = 3$pp, $p = 0.31$, Cohen's $d = 0.12$).
At 28\% base success, the agent's ``correction'' is statistically indistinguishable from random re-sampling.

\subsubsection{B-IAVR Training}

\begin{table}[h]
\centering
\caption{B-IAVR selective verification training results.}
\label{tab:biavr}
\begin{tabular}{lccc}
\toprule
Metric & Target & Achieved & Gap \\
\midrule
Verification rate reduction & $\geq$40\% & 12\%   & $-$28pp \\
Success $\Delta$ vs. always-on & $<$5pp   & $-$6.2pp & Fail \\
Cost reduction               & $\geq$40\% & 12\%   & $-$28pp \\
Policy entropy (final)       & $>$0.3    & 0.08   & Collapsed \\
\bottomrule
\end{tabular}
\end{table}

The B-IAVR policy degenerates to near-always-verify (88\% verification rate, Table~\ref{tab:biavr}).
The policy entropy collapses to 0.08, indicating the policy learned a trivial solution: always verify with occasional random skips.
Success drops 6.2pp below always-on verification, and cost reduction is only 12\%, far below the 40\% target.
The root cause is twofold: (1) the action-change reward signal is noise (Section 4.6.1), providing no learning gradient; (2) at 28\% base success, any skip risks the already-fragile performance, making always-verify the locally optimal degenerate policy.

\subsubsection{Verification Strategy Comparison}

\begin{table}[h]
\centering
\caption{Verification strategy comparison across methods.}
\label{tab:verifycomp}
\begin{tabular}{lcccc}
\toprule
Method & Success & Verify Rate & Relative Cost & Success/Token \\
\midrule
Always-on         & 28.1\% & 100\% & 1.00$\times$ & 0.281 \\
B-IAVR            & 21.9\% & 88\%  & 0.92$\times$ & 0.238 \\
Random ($p{=}0.5$) & 23.4\% & 50\%  & 0.67$\times$ & 0.349 \\
Confidence        & 25.0\% & 62\%  & 0.74$\times$ & 0.338 \\
No verify         & 18.8\% & 0\%   & 0.45$\times$ & 0.418 \\
\bottomrule
\end{tabular}
\end{table}

Table~\ref{tab:verifycomp} shows B-IAVR is Pareto-dominated by every baseline.
Random verification at $p{=}0.5$ achieves comparable success (23.4\% vs. 21.9\%) at nearly half the verification cost.
The confidence-based heuristic (25.0\% at 62\% verification) also dominates B-IAVR.
The no-verify condition drops only 9.3pp from always-on (28.1\% to 18.8\%), indicating that verification itself has limited value at this performance level: the agent's corrections are near-random perturbations.

\subsubsection{Reward Signal Ablation}

\begin{table}[h]
\centering
\caption{Reward signal comparison for selective verification training.}
\label{tab:rewardablation}
\begin{tabular}{lccc}
\toprule
Reward Type & Success & Verify Rate & Pareto-Dominated? \\
\midrule
Action-change  & 21.9\% & 88\% & Yes (by all others) \\
Random trigger & 23.4\% & 52\% & No \\
Confidence     & 25.0\% & 65\% & No \\
\bottomrule
\end{tabular}
\end{table}

The action-change reward, our proposed intrinsic signal, produces the worst selective policy (Table~\ref{tab:rewardablation}).
It trains a near-always-verify policy because the reward is noise (partial correlation with success: 0.03 after entropy control).
Random triggering learns a more balanced skip rate because uniform random reward at least does not actively mislead.
Confidence-based reward performs best, as model confidence has weak but real correlation with success.

\subsubsection{ARES Comparison}

ARES, which uses supervised learning from oracle verification decisions, achieves 25.0\% success with 55\% verification rate and 38\% cost reduction.
B-IAVR achieves 21.9\% at 88\% verification rate with 12\% cost reduction.
ARES dominates B-IAVR on all metrics.
This is expected: supervised learning from oracle labels does not require a meaningful RL reward signal, bypassing the core bottleneck that causes B-IAVR to fail.
The result suggests that SFT-based verification routing is more practical than RL-based approaches when base agent competence is low.

\subsection{Transfer and Scaling}

\subsubsection{ScienceWorld Transfer}

\begin{table}[h]
\centering
\caption{ScienceWorld transfer: episode-level vs. step-independent.}
\label{tab:scienceworld}
\begin{tabular}{lccc}
\toprule
Method & Success & Std & Avg Episode Length \\
\midrule
Episode-level         & 15.0\% & 6pp & 12.3 steps \\
Step-independent ($h{=}5$) & 12.0\% & 8pp & 12.3 steps \\
\bottomrule
\end{tabular}
\end{table}

On ScienceWorld (Table~\ref{tab:scienceworld}), episode-level GRPO achieves only a 3pp advantage over step-independent ($p = 0.38$, not significant).
ScienceWorld episodes average 12.3 steps, compared to 15--45 in ALFWorld.
With shorter episodes, step-independent windows capture a larger fraction of the trajectory, reducing information loss and sign misattribution.
This result indicates that the step-independent disadvantage is \emph{horizon-dependent}: it manifests primarily in long-episode environments.

\subsubsection{7B Model Scaling}

\begin{table}[h]
\centering
\caption{Scaling to Qwen3-7B: episode-level vs. step-independent.}
\label{tab:scaling}
\begin{tabular}{lccc}
\toprule
Method & 1.7B & 7B & Gap Reduction \\
\midrule
Episode-level            & 28.1\% & 35\% & -- \\
Step-independent ($h{=}5$) & 0\%    & 22\% & 28pp $\to$ 13pp \\
\bottomrule
\end{tabular}
\end{table}

At 7B scale (Table~\ref{tab:scaling}), step-independent training no longer collapses: it reaches 22\% success at $h{=}5$, compared to 0\% at 1.7B.
The gap between episode-level and step-independent narrows from 28pp to 13pp.
Larger models have stronger in-context learning capabilities, partially compensating for the limited context window in step-independent training.
This suggests that the ``catastrophic failure'' of step-independent training is specific to smaller models on long-horizon tasks.

\subsubsection{Per-Task Breakdown}

\begin{table}[h]
\centering
\caption{Per-task-type breakdown (episode-level GRPO, 1.7B).}
\label{tab:pertask}
\begin{tabular}{lccc}
\toprule
Task Type & Episode-Level & Step-Indep.\ ($h{=}5$) & Avg Length \\
\midrule
Pick       & 35\% & 8\%  & 15 steps \\
Heat       & 30\% & 3\%  & 22 steps \\
Clean      & 25\% & 0\%  & 28 steps \\
Cool       & 20\% & 2\%  & 25 steps \\
Examine    & 15\% & 0\%  & 38 steps \\
Pick\_two  & 10\% & 0\%  & 45 steps \\
\bottomrule
\end{tabular}
\end{table}

Table~\ref{tab:pertask} shows that both episode-level success and the episode-level vs. step-independent gap correlate with task length (correlation 0.45 with episode length).
Shorter tasks (Pick, 15 steps) are easiest; longer multi-object tasks (Pick\_two, 45 steps) are hardest.
Step-independent training fails completely on tasks longer than 25 steps.
However, the moderate correlation (0.45) indicates that task complexity, not length alone, drives difficulty.


%% ============================================================
%% 5  DISCUSSION
%% ============================================================
\section{Discussion}

\paragraph{The coarser-is-sufficient hypothesis is false.}
Our original hypothesis was that episode-level credit assignment, while coarser than per-token methods, would suffice for training effective multi-turn agents.
The 20pp gap between episode-level GRPO and GAE($\lambda{=}0.95$) definitively refutes this.
Episode-level training provides correct advantage direction (success vs. failure) but cannot differentiate token-level contributions, leading to an early plateau.

\paragraph{Sign misattribution is real but not dominant.}
We find empirical support for the advantage sign misattribution mechanism (0.72 sign-outcome correlation), but it is weaker than theoretically predicted.
GRPO's group normalization introduces partial local ranking that our theory did not account for, resulting in 18\% mixed-sign episodes and a non-trivial 0.28 correlation with local action quality.
Sign misattribution is one factor in step-independent failure, alongside information loss and $\gamma$-decay, rather than the sole mechanism.

\paragraph{Selective verification has a competence prerequisite.}
The complete failure of B-IAVR ($-$6.2pp success, 12\% cost reduction, Pareto-dominated by all baselines) reveals a prerequisite not discussed in prior work: the base agent must be competent enough for verification to have causal value.
At 28\% success, verification-triggered action changes are statistically indistinguishable from random perturbations (3pp difference, $p = 0.31$).
We conjecture that a base agent competence of $\geq$50\% is necessary for selective verification to be learnable via RL, though we cannot test this with our training setup.

\paragraph{Implications for the field.}
Our results suggest two actionable takeaways:
(1) Practitioners should not assume step-independent training will work for long-horizon agent tasks, particularly at small model scales, but the solution is finer credit (GAE or similar) rather than coarser episode-level training.
(2) Selective verification research should report base agent competence as a key variable and establish minimum competence thresholds before pursuing cost-efficiency optimizations.


%% ============================================================
%% 6  CONCLUSION
%% ============================================================
\section{Conclusion}

We presented an empirical taxonomy of credit assignment failure modes in multi-turn agent RL, characterized across granularity, agent competence, and task horizon.
Step-independent training exhibits a stochastic, seed-dependent performance cliff that is most severe at small model scales and long task horizons.
Episode-level GRPO avoids collapse but plateaus significantly below GAE-based training, demonstrating that per-token credit differentiation is necessary.
Selective self-verification fails at base agent competence below 30\%, revealing a prerequisite bottleneck.
These negative results constrain the design space for future work: neither ``coarser credit suffices'' nor ``learn when to verify on a weak agent'' are viable research directions, while the interaction between model scale, task horizon, and credit granularity merits systematic investigation.

\subsection*{Limitations}

Our study is limited to ALFWorld and ScienceWorld with models up to 7B parameters.
The step-independent failure may not generalize to environments with denser reward signals or shorter episodes.
The multi-seed analysis uses 3--5 seeds, which provides limited statistical power for characterizing distributional properties.
The GAE results use a simple value function architecture; more sophisticated critics may further widen the gap.
We could not test the conjectured 50\% competence threshold for selective verification because our training setup does not reach that level.

\bibliographystyle{plainnat}
\bibliography{references}

\end{document}
